\section{全讲总结}

这一讲其实只想回答一个问题：\textbf{当模型规模和硬件规模一起变大时，训练系统应该怎样分配计算与通信？}

\subsection{本章小结}

\begin{importantbox}{三条最终 takeaway}
\begin{enumerate}
    \item \textbf{数据并行} 解决的是 batch 维度扩展，最简单也最常用。
    \item \textbf{张量并行} 解决的是单层太宽的问题，但它非常吃互联带宽。
    \item \textbf{流水线并行} 解决的是模型太深的问题，但它要面对 bubble 和调度复杂度。
\end{enumerate}
\end{importantbox}

\begin{table}[H]
\centering
\small
\begin{tabular}{p{3.0cm}p{3.5cm}p{3.0cm}p{3.5cm}}
\toprule
\textbf{核心概念} & \textbf{一句话解释} & \textbf{关键公式} & \textbf{工程要点} \\
\midrule
集合通信 & 多 rank 间的标准协议 & $T = \alpha + S/B$ & NCCL 自动优化路径 \\
Ring all-reduce & 带宽最优的全局归约 & $2\frac{p-1}{p}S$ & 大张量时接近理论峰值 \\
数据并行 & 复制模型，切分 batch & 每步 all-reduce $O(P)$ & 先从 DDP 开始 \\
张量并行 & 切分层宽度 & 每层 all-gather $O(A)$ & 仅限 NVLink 范围 \\
流水线并行 & 切分层深度 & bubble $\frac{S-1}{M+S-1}$ & microbatch 数要够大 \\
ZeRO & 分片训练状态 & 每卡 $\frac{16P}{p}$ & ZeRO-2 性价比最高 \\
\bottomrule
\end{tabular}
\caption{全讲核心概念速查表}
\end{table}

\subsection{这节课没有覆盖的内容}

讲者在最后提到了几个本讲没有深入的主题：

\begin{itemize}
    \item \textbf{更通用的模型}：本讲用 MLP 演示，真实 transformer 还有 attention 层需要特殊的分片策略（如 Megatron 的 column/row parallel）。
    \item \textbf{通信/计算重叠}：DDP 的 bucket gradient all-reduce、1F1B 调度等技术可以让通信隐藏在计算背后，但代码实现复杂度会大幅上升。
    \item \textbf{序列并行}（Sequence Parallelism）：当序列长度非常长时，可以沿 sequence 维度切分——这是除了 batch、width、depth 之外的第四个维度。
    \item \textbf{JAX/TPU 的声明式方法}：在 JAX 生态中，用户只需声明分片策略，编译器自动插入通信原语。PyTorch 则要求手动编排，但能看清底层发生了什么。
\end{itemize}

\begin{knowledgebox}{最后一句话}
如果你记住一件事，那就是：多 GPU 训练的本质不是让 GPU 彼此热闹起来，而是让通信次数、通信量和同步频率，恰好配得上模型规模。
\end{knowledgebox}

